\documentclass[10pt,a4paper]{article}

\usepackage[T1]{fontenc}
\usepackage[utf8]{inputenc}
\usepackage[ngerman]{babel}
\usepackage[a4paper,margin=14mm]{geometry}
\usepackage{lmodern}
\usepackage[table]{xcolor}
\usepackage{tabularx}
\usepackage{array}
\usepackage{microtype}
\usepackage[most]{tcolorbox}

\pagestyle{empty}
\setlength{\parindent}{0pt}
\setlength{\parskip}{0pt}
\renewcommand{\familydefault}{\sfdefault}

\definecolor{Navy}{HTML}{17324D}
\definecolor{Blue}{HTML}{3B6EA5}
\definecolor{Sky}{HTML}{EAF3FA}
\definecolor{Mint}{HTML}{E8F5EF}
\definecolor{Gold}{HTML}{F3B33D}
\definecolor{Ink}{HTML}{1E2935}
\definecolor{Muted}{HTML}{617080}
\definecolor{Line}{HTML}{D8E1E8}
\definecolor{CodeBg}{HTML}{F5F7F9}
\definecolor{Cell}{HTML}{FFF4DB}
\definecolor{GridHead}{HTML}{E9EDF1}
\definecolor{GridLine}{HTML}{AEBBC6}

\color{Ink}

\newcommand{\sectiontitle}[1]{%
  \vspace{2.3mm}
  {\color{Blue}\bfseries\large #1}\par
  \vspace{0.8mm}
  {\color{Blue}\rule{\linewidth}{0.7pt}}\par
  \vspace{1.2mm}
}

% Tabellenausschnitt: Kopfzeile (Spaltenbuchstaben) und Zeilennummern
\newcommand{\gridhead}[1]{\cellcolor{GridHead}\textcolor{Muted}{\bfseries\small #1}}
\newcommand{\gridrow}[1]{\cellcolor{GridHead}\textcolor{Muted}{\bfseries\small #1}}
\newcommand{\sheetsetup}{%
  \arrayrulecolor{GridLine}\setlength{\arrayrulewidth}{0.4pt}%
  \renewcommand{\arraystretch}{1.25}}

\begin{document}

\colorbox{Navy}{%
  \parbox{\dimexpr\linewidth-2\fboxsep\relax}{%
    \vspace{3mm}\color{white}
    \begin{tabularx}{\linewidth}{@{}Xr@{}}
      {\small\bfseries INFORMATIK · KLASSE 9} &
      \colorbox{Gold}{\color{Navy}\small\bfseries\strut LÖSUNGEN}
    \end{tabularx}
    \vspace{2.5mm}

    {\fontsize{22}{25}\selectfont\bfseries Sicherungsblatt: Aufgabe 1}\par
    \vspace{1mm}
    {\large Formeln in der Tabellenkalkulation}
    \vspace{3mm}
  }%
}

\vspace{3mm}
\colorbox{Sky}{%
  \parbox{\dimexpr\linewidth-2\fboxsep\relax}{%
    \vspace{1.8mm}
    \textcolor{Blue}{\bfseries Ziel:}
    Du benennst Zellen und rechnest mit Formeln und Zellbezügen.
    \vspace{1.8mm}
  }%
}

\sectiontitle{1 · Zellen benennen}

\begin{minipage}[t]{62mm}
  \vspace{0pt}
  \sheetsetup
  \begin{tabular}{|c|c|c|c|}
    \hline
    \cellcolor{GridHead} & \gridhead{A} & \gridhead{B} & \gridhead{C} \\ \hline
    \gridrow{3} & \hspace{13mm} & \hspace{13mm} & \hspace{13mm} \\ \hline
    \gridrow{4} & & & \\ \hline
    \gridrow{5} & & \cellcolor{Cell}\small\bfseries B5 & \\ \hline
    \gridrow{6} & & & \\ \hline
  \end{tabular}
\end{minipage}\hfill
\begin{minipage}[t]{\dimexpr\linewidth-68mm\relax}
  \vspace{2mm}
  Spalte \textbf{B}, Zeile \textbf{5} $\Rightarrow$ Zelle \textbf{B5}\par
  \vspace{2mm}
  Programme: Microsoft Excel, LibreOffice Calc
\end{minipage}

\sectiontitle{2 · Formeln}

Jede Formel beginnt mit \textbf{=}.

\vspace{1.5mm}
\renewcommand{\arraystretch}{1.2}
\begin{tabularx}{\linewidth}{@{}>{\centering\ttfamily\bfseries}p{9mm}X>{\centering\ttfamily\bfseries}p{9mm}X>{\centering\ttfamily\bfseries}p{9mm}X@{}}
  \rowcolor{CodeBg} + & plus & * & mal & \^{} & hoch \\
  - & minus & / & geteilt durch & ( ) & Klammern \\
\end{tabularx}

\vspace{2mm}
\begin{tabularx}{\linewidth}{@{}X X@{}}
  \colorbox{CodeBg}{\parbox{\dimexpr\linewidth-4mm\relax}{%
    \texttt{=3*4} \quad{\small liefert immer 12}}}
  &
  \colorbox{Mint}{\parbox{\dimexpr\linewidth-4mm\relax}{%
    \texttt{=B4*B5} \quad{\small rechnet mit neuen Werten neu}}}
  \\
\end{tabularx}

\sectiontitle{3 · Ausschnitt aus Aufgabe 1}

\sheetsetup
\begin{tabular}{|c|l|c|c|l|l|}
  \hline
  \cellcolor{GridHead} & \gridhead{A} & \gridhead{B} & \gridhead{\dots} &
  \gridhead{E} & \gridhead{F} \\ \hline
  \gridrow{4} & \small Seite a / Grundseite g: & \small 3 & &
    \small A\_Dreieck = & \cellcolor{Cell}\ttfamily\small =B4*B8/2 \\ \hline
  \gridrow{5} & \small Seite b: & \small 4 & & \small A\_Rechteck = & \\ \hline
  \gridrow{6} & \small Seite c: & \small 5 & &
    \small A\_Quadrat = & \cellcolor{Cell}\ttfamily\small =B4\^{}2 \\ \hline
  \gridrow{7} & \small Seite d / Körperhöhe: & \small 6 & & \small A\_Trapez = & \\ \hline
  \gridrow{8} & \small Höhe h: & \small 7 & & \small A\_Kreis = & \\ \hline
\end{tabular}

\vspace{1.5mm}
{\small In der Zelle stehen die Ergebnisse \textbf{10,5} und \textbf{9},
in der Formelleiste die Formel.\par}

\sectiontitle{4 · Merke}

\begin{tcolorbox}[colback=Mint,colframe=Blue!40,arc=3mm,boxrule=0.7pt,
  left=3mm,right=3mm,top=2mm,bottom=2mm]
  Eine Zelle wird durch Spalte und Zeile bezeichnet: \textbf{B5}.\par
  \vspace{1.2mm}
  Jede Formel beginnt mit \texttt{=} und rechnet mit
  \texttt{+ - * / \^{}} und Klammern.\par
  \vspace{1.2mm}
  Formeln mit Zellbezügen rechnen bei neuen Werten automatisch neu.
\end{tcolorbox}

\vfill
{\color{Line}\rule{\linewidth}{0.5pt}}\par
\vspace{1.5mm}
\begin{tabularx}{\linewidth}{@{}Xr@{}}
  {\small\color{Muted}Klasse 9 · Tabellenkalkulation} &
  {\small\color{Muted}Sicherungsblatt · Aufgabe 1 · Version 2}
\end{tabularx}

\end{document}
